\documentclass{article}
\usepackage{amsmath, amssymb}
\usepackage{geometry}
\geometry{margin=1in}
\begin{document}

\title{CSE 400: Fundamentals of Probability in Computing\\
Tutorial 2 – Lecture Scribe (Step-by-Step)}
\date{}
\maketitle

\section*{Q1. Random Point on a Line Segment}

\textbf{Problem Interpretation}

A point is chosen uniformly at random on a line segment of length $L$.

Let:

The point be at distance $x$ from one end.

Then $0 < x < L$.

The two segments formed are:

$x$

$L - x$

The ratio of the shorter to the longer segment is:

\[
R = \frac{\min(x, L-x)}{\max(x, L-x)}
\]

We must find:

\[
P(R < 1/4)
\]

\textbf{Logical Case Division}

We divide based on midpoint $L/2$.

\textbf{Case 1:} $x \le L/2$

Shorter segment = $x$  
Longer segment = $L - x$

\[
\frac{x}{L-x} < \frac{1}{4}
\]

Multiply both sides:

\[
4x < L - x
\]

\[
5x < L
\]

\[
x < \frac{L}{5}
\]

\textbf{Case 2:} $x > L/2$

Shorter segment = $L - x$  
Longer segment = $x$

\[
\frac{L-x}{x} < \frac{1}{4}
\]

Multiply:

\[
4L - 4x < x
\]

\[
5x > 4L
\]

\[
x > \frac{4L}{5}
\]

\textbf{Valid Region}

\[
(0, L/5) \cup (4L/5, L)
\]

Total length:

\[
(L/5 - 0) + (L - 4L/5) = L/5 + L/5 = 2L/5
\]

Since uniform over length $L$:

\[
P = \frac{2L/5}{L} = \frac{2}{5}
\]

\section*{Q2. Bayes Classification with Normal Densities}

\textbf{Given}

White region:

\[
X|W \sim N(4,4), \sigma_1 = 2
\]

Black region:

\[
X|B \sim N(6,9), \sigma_2 = 3
\]

Prior probabilities:

\[
P(B) = \alpha, \quad P(W) = 1-\alpha
\]

Observed:

\[
X = 5
\]

We want equal posterior error:

\[
P(B|X=5) = P(W|X=5) = 1/2
\]

\textbf{Bayes Formula}

\[
P(B|X=5) = \frac{\alpha f(5|B)}{\alpha f(5|B) + (1-\alpha)f(5|W)}
\]

Set equal to $1/2$.

\textbf{Normal Density Formula}

\[
f(x) = \frac{1}{\sigma \sqrt{2\pi}} e^{-\frac{(x-\mu)^2}{2\sigma^2}}
\]

Compute densities:

\[
f(5|B) = \frac{1}{3\sqrt{2\pi}} e^{-1/18}
\]

\[
f(5|W) = \frac{1}{2\sqrt{2\pi}} e^{-1/8}
\]

Substitute:

\[
\frac{\alpha e^{-1/18}/(3\sqrt{2\pi})}
{\alpha e^{-1/18}/(3\sqrt{2\pi}) + (1-\alpha)e^{-1/8}/(2\sqrt{2\pi})}
= \frac{1}{2}
\]

Simplifies to:

\[
\frac{3\alpha}{3\alpha + 2(1-\alpha)e^{-5/72}} = \frac{1}{2}
\]

Solve:

\[
6\alpha = 3\alpha + 2(1-\alpha)(0.93)
\]

\[
3\alpha = 1.86 - 1.86\alpha
\]

\[
4.86\alpha = 1.86
\]

\[
\alpha \approx 0.3827
\]

\section*{Q3. Fire Station Location}

\textbf{(a) Uniform(0, A)}

\[
X \sim Uniform(0,A)
\]

\[
f_X(x) = \frac{1}{A}
\]

Minimize:

\[
E[|X-a|]
\]

Split integral:

\[
\int_0^a (a-x)\frac{1}{A}dx + \int_a^A (x-a)\frac{1}{A}dx
\]

Differentiate:

\[
\frac{d}{da}\left(\frac{2a^2 - 2aA + A^2}{2A}\right)
= \frac{1}{2A}(4a - 2A)
\]

Set zero:

\[
4a = 2A
\]

\[
a = \frac{A}{2}
\]

Station at midpoint.

\textbf{(b) Exponential($\lambda$)}

\[
f_X(x) = \lambda e^{-\lambda x}
\]

Differentiate expectation:

\[
\frac{d}{da}E[|X-a|] = 1 - 2e^{-\lambda a}
\]

Set zero:

\[
2e^{-\lambda a} = 1
\]

\[
e^{-\lambda a} = 1/2
\]

\[
-\lambda a = \ln(1/2) = -\ln 2
\]

\[
a = \frac{\ln 2}{\lambda}
\]

\section*{Q4. Mixture of Exponentials}

Two battery types:

\[
P(T=i) = p_i
\]

\[
P(X>t) = p_1 e^{-\lambda_1 t} + p_2 e^{-\lambda_2 t}
\]

Similarly:

\[
P(X>t+s) = p_1 e^{-\lambda_1 (t+s)} + p_2 e^{-\lambda_2 (t+s)}
\]

Conditional probability:

\[
P(X>t+s|X>t)
=
\frac{p_1 e^{-\lambda_1 (t+s)} + p_2 e^{-\lambda_2 (t+s)}}
{p_1 e^{-\lambda_1 t} + p_2 e^{-\lambda_2 t}}
\]

\section*{Q5. Conditional Poisson}

\[
X \sim Poisson(\lambda_1)
\]

\[
Y \sim Poisson(\lambda_2)
\]

Independent.

Find:

\[
P(X=k|X+Y=n)
\]

Step 1:

\[
P(X=k, Y=n-k) = P(X=k)P(Y=n-k)
\]

Substitute Poisson PMF:

\[
= \frac{e^{-\lambda_1}\lambda_1^k}{k!}
\cdot
\frac{e^{-\lambda_2}\lambda_2^{n-k}}{(n-k)!}
\]

Denominator:

\[
P(X+Y=n)
=
\frac{e^{-(\lambda_1+\lambda_2)}(\lambda_1+\lambda_2)^n}{n!}
\]

Final Simplified Form:

\[
P(X=k|X+Y=n)
=
\binom{n}{k}
\left(\frac{\lambda_1}{\lambda_1+\lambda_2}\right)^k
\left(\frac{\lambda_2}{\lambda_1+\lambda_2}\right)^{n-k}
\]

Binomial structure appears.

\section*{Q6. Transformation of Random Variable}

\[
X \sim Uniform[-2,2]
\]

\[
f_X(x) = \frac{1}{4}
\]

Piecewise function:

$x$ for $[-2,-1]$

$0$ for $(-1,1)$

$x$ for $[1,2]$

CDF

Compute piecewise:

$y<-2$: 0

$-2 \le y < -1$: $(y+2)/4$

$-1 \le y < 0$: $1/4$

At $y=0$: jump of $1/2$

\[
F_Y(0) = 1/4 + 1/2 = 3/4
\]

$0<y<1$: $3/4$

$1 \le y \le 2$: $(y+2)/4$

$y>2$: 1

PDF

\[
f_Y(y) = 1/4
\]
for $(-2,-1)$ and $(1,2)$

Point mass:

\[
P(Y=0) = 1/2
\]

\section*{Q7. Stock Model}

\[
S_{1000} = S_0 u^k d^{1000-k}
\]

Require:

\[
u^k d^{1000-k} \ge 1.30
\]

Take logs:

\[
k \ln u + (1000-k)\ln d \ge \ln(1.30)
\]

Solve:

\[
k \ge
\frac{\ln(1.30) - 1000\ln d}
{\ln u - \ln d}
\]

Substitute values:

\[
k \ge 469.4
\]

Thus:

\[
k \ge 470
\]

Since:

\[
k \sim Binomial(1000,0.52)
\]

Probability:

\[
\sum_{k=470}^{1000}
\binom{1000}{k}
(0.52)^k
(0.48)^{1000-k}
\]

\section*{Q8. Exponential Distribution}

\[
X \sim Exponential(1/2)
\]

PDF:

\[
f(x) = \frac{1}{2} e^{-x/2}
\]

(a)

\[
P(X>2) = e^{-1} \approx 0.3679
\]

(b)

\[
P(X \ge 10 | X>9)
=
\frac{e^{-5}}{e^{-9/2}}
=
e^{-1/2}
\]

\section*{Q9. Distribution of $Z = X^2 + Y^2$}

\[
F_Z(z) = P(X^2 + Y^2 \le z^2)
\]

Region:

\[
x^2 + y^2 \le z^2
\]

\[
F_Z(z)
=
\int_{-z}^{z}
\int_{-\sqrt{z^2-x^2}}^{\sqrt{z^2-x^2}}
f_X(x)f_Y(y)
\, dy \, dx
\]

Differentiate using Leibniz rule:

\[
\frac{d}{dz}\sqrt{z^2 - x^2}
=
\frac{z}{\sqrt{z^2-x^2}}
\]

Final PDF:

\[
f_Z(z)
=
z
\int_{-z}^{z}
f_X(x)
\frac{
f_Y(\sqrt{z^2-x^2}) + f_Y(-\sqrt{z^2-x^2})
}
{\sqrt{z^2-x^2}}
dx
\]

\section*{Q10. Joint Distribution of $M_n, M_{n+1}$}

\[
M_n = \max(X_1,\dots,X_n)
\]

\[
M_{n+1} = \max(M_n, X_{n+1})
\]

CDF of $M_n$

\[
P(M_n \le x) = F(x)^n
\]

Case 1: $x \le y$

\[
P(M_n \le x, M_{n+1} \le y)
=
F(x)^n F(y)
\]

Case 2: $x > y$

\[
P(M_n \le x, M_{n+1} \le y)
=
F(y)^{n+1}
\]

Final Joint CDF:

\[
P(M_n \le x, M_{n+1} \le y)
=
\begin{cases}
F(x)^n F(y), & x \le y \\
F(y)^{n+1}, & x > y
\end{cases}
\]

\end{document}